% Part: lambda-calculus
% Chapter: introduction
% Section: composition

\documentclass[../../../include/open-logic-section]{subfiles}

\begin{document}

\olfileid[fr]{lam}{int}{com}
\olsection{Les fonctions \texorpdfstring{$\lambd$}{λ}-définissables sont fermées par composition}

\begin{lem}
Les fonctions totales $\lambd$-définissables sont
fermées par composition.
\end{lem}

\begin{proof}
Supposons que $f$ soit définie par composition
à partir de $h$, $g_0$, \dots,~$g_{k-1}$.
Si $h$, $g_0$, \dots,~$g_{k-1}$ sont respectivement
définies par les termes $H$, $G_0$, \dots,~$G_{k-1}$,
nous cherchons un terme~$F$ qui définisse~$f$.
Posons
\[
F(x_0, \dots, x_{l-1}) = H(G_0(x_0, \dots, x_{l-1}),
\dots, G_{k-1}(x_0, \dots, x_{l-1})).
\]
Lorsque toutes les valeurs intermédiaires existent,
ce terme donne bien la valeur de la composition.
La notation du lambda-calcul se prête ainsi
à la représentation de la composition.
\end{proof}

Pour des fonctions partielles, la composition exige que tous les
arguments de la fonction extérieure soient définis. Posons
$\Omega\ident(\lambd[z][zz])(\lambd[z][zz])$,
$H(v)\ident\num{0}$ et $G_0(u)\ident\Omega$.
La composition stricte de ces fonctions est indéfinie partout,
tandis que $H(G_0(x))$ se réduit à~$\num{0}$. La fermeture de la
classe des fonctions partielles $\lambd$-définissables découle
des deux théorèmes de représentation
(\olref[lam][rep][par]{thm:partial-recursive-lambda} et
\olref[lam][rep][ldr]{thm:lambda-computable}) et de la clôture
des fonctions récursives partielles par composition.

\end{document}
